% TOP_PROBLEM: {"id": 3059, "title": "A nowhere-zero point in a linear mapping", "category_id": 2, "category": {"id": 2, "name": "combinatorics", "display_name": "Combinatorics", "description": "Counting problems, graph theory, discrete structures.", "slug": "combinatorics", "order_index": 2, "created_at": "2026-07-31T15:26:25.671Z"}, "status": "open", "problem_number": "OPG-148", "set_id": 12}
% FIRST_POSTED: 2026-10-01
% ATTEMPT_STATUS: unresolved
% UnsolvedMath ID 3059; source catalogue code OPG-148
% !TeX program = lualatex
% GPT-6 Astra Ultra research attempt; independent partial work, 2026-10-01.
\documentclass[11pt,a4paper]{article}
\usepackage[margin=25mm]{geometry}
\usepackage{fontspec}
\setmainfont{lmroman10-regular.otf}[BoldFont=lmroman10-bold.otf,
 ItalicFont=lmroman10-italic.otf,BoldItalicFont=lmroman10-bolditalic.otf]
\usepackage{amsmath,amssymb,mathtools}
\usepackage{unicode-math}
\setmathfont{latinmodern-math.otf}
\AtBeginDocument{\let\setminus\smallsetminus}
\usepackage{xurl,hyperref}
\hypersetup{colorlinks=true,urlcolor=blue}
\setcounter{secnumdepth}{0}
\setlength{\parindent}{0pt}
\setlength{\parskip}{5pt}
\setlength{\emergencystretch}{3em}
\begin{document}
\section*{A nowhere-zero point in a linear mapping}\label{3059}
{\small UnsolvedMath ID 3059 · OPG-148 · Combinatorics · OpenGarden}
\subsection{Definitions and mathematical statement}
Let $\mathbb F_q$ be a finite field of cardinality $q\ge4$, let $n\ge1$,
and let $A\in\mathbb F_q^{n\times n}$ be invertible, meaning
$\det(A)\ne0$ in the field. Put
$\mathbb F_q^\times=\mathbb F_q\setminus\{0\}$ and
$T=(\mathbb F_q^\times)^n$.
A vector is nowhere zero when each of its coordinates is nonzero; this
is stronger than requiring the vector itself to differ from the zero vector.

The conjecture states
\[
 \forall q\ge4\text{ a prime power}\quad\forall n\ge1\quad
 \forall A\in\operatorname{GL}_n(\mathbb F_q),\qquad
 A(T)\cap T\ne\varnothing.
\]
Equivalently, there are $x_1,\ldots,x_n\in\mathbb F_q^\times$ such that
$\sum_{j=1}^n a_{ij}x_j\ne0$ for every row $i$. Setting $y=Ax$ then
supplies the two nowhere-zero vectors appearing in the original statement.
Both vectors are over the same field; their coordinates may repeat.

Another description asks whether the $n$ coordinate hyperplanes
$x_j=0$, together with the $n$ hyperplanes
$\sum_j a_{ij}x_j=0$, can ever cover all of $\mathbb F_q^n$ when the rows
of $A$ are independent. The desired answer is that they cannot.
The matrix and dimension are arbitrary, and any existence argument must
work without additional assumptions on the pattern of nonzero entries.
The lower bound $q\ge4$ is part of the assertion and is retained.
\subsection{Short English statement}
Must every invertible matrix over a finite field with at least four elements send some coordinatewise nonzero vector to another such vector?
\subsection{Research attempt: exact row probabilities and a triangular subcase}
\paragraph{Scope and literature check.}
GPT-6 Astra Ultra, 1 October 2026. The coordinates must all be nonzero;
avoiding only the zero vector would answer a different question. Alon
and Tarsi [S3] prove the assertion for nonprime finite fields. Thus those
fields are already covered by literature. This attempt develops direct
elementary certificates, valid in every characteristic in their stated
ranges, and does not reproduce the full polynomial theorem.

\paragraph{1. Count a single forbidden row exactly.}
Choose $X$ uniformly from $(\mathbb F_q^\times)^n$. A row with $s$
nonzero coefficients has the same zero probability as
$Z_1+\cdots+Z_s=0$, with independent uniform nonzero $Z_i$; multiplying
a coordinate by a nonzero coefficient preserves its distribution.
Let $N_s$ be the number of zero-sum $s$-tuples. We have $N_1=0$ and
\[
 N_s=(q-1)^{s-1}-N_{s-1}.
\]
Indeed the final entry is uniquely determined and nonzero exactly when
the sum of the preceding entries is nonzero. Solving the recurrence gives
\[
 N_s=\frac{(q-1)^s+(q-1)(-1)^s}{q},\qquad
 p_s=\frac1q\left(1+\frac{(-1)^s}{(q-1)^{s-1}}\right).
\]
In particular $p_1=0$ and $p_s\le1/(q-1)$ for $s\ge2$.
If the row support sizes are $s_1,\ldots,s_n$, the union bound yields
\[
 |\{x\in(\mathbb F_q^\times)^n:Ax\in(\mathbb F_q^\times)^n\}|
 \ge(q-1)^n\left(1-\sum_i p_{s_i}\right).
\]
Thus $\sum_i p_{s_i}<1$ is a checkable sufficient condition. In particular,
the conjecture holds for $n<q-1$, and more generally if fewer than $q-1$
rows have support size at least two. Singleton rows impose no forbidden
event at all. Invertibility guarantees every row is nonzero, but this
union-bound proof otherwise does not exploit independence of the rows.

\paragraph{2. Arbitrary dimension for triangular matrices.}
Suppose $A$ is upper triangular and has nonzero diagonal. Choose the
coordinates in order $x_n,x_{n-1},\ldots,x_1$. Once $x_{i+1},\ldots,x_n$
are chosen, row $i$ forbids at most one nonzero value of $x_i$, namely
the possible solution of
$a_{ii}x_i=-\sum_{j>i}a_{ij}x_j$. Since there are $q-1$ nonzero
choices, at least $q-2$ choices remain when $q>2$.
Earlier chosen output coordinates cannot change, since rows below $i$
do not involve $x_i$. This constructs a successful vector, and even
at least $(q-2)^n$ successful vectors. The result also holds after row
or column permutations and nonzero coordinate scalings, since these
operations preserve the coordinate torus.

\paragraph{3. Why Gaussian elimination is not a solution.}
General row addition does not preserve the output torus: in any field,
two nonzero coordinates can sum to zero. Thus reducing $A$ to triangular
form by arbitrary row operations does not transfer the property back.
The small-field matrix
\[
 A=\begin{pmatrix}1&1\\1&-1\end{pmatrix}\quad\hbox{over }\mathbb F_3
\]
is invertible but has no successful torus vector: for nonzero $x_1,x_2$,
either $x_2=x_1$ or $x_2=-x_1$, and one output coordinate vanishes.
This also checks the essential field-size restriction.

\paragraph{Verification and remaining gap.}
The recurrence gives $p_2=1/(q-1)$ and
$p_3=(q-2)/(q-1)^2$, agreeing with direct enumeration. Dependencies
between row events were never treated as probabilistic independence.
For large dense matrices over prime fields the sum of the row
probabilities can exceed one, giving no conclusion. The first missing
ingredient is a way to use linear independence of the rows to control
their overlapping zero sets beyond a union bound.

\paragraph{Reproduction of finite checks.}
The finite checks stated above were run with Python 3 using the repository
script [C1]. They verify the stated examples and finite ranges only.

\subsection{Final status}
UNRESOLVED as a general independent attempt. The support criterion and
triangular case are proved; the literature already settles nonprime
fields. No new solution or quantitative completion percentage is claimed.

\subsection{Sources}
\begin{itemize}
\item[S1] ULAM.AI and the UnsolvedMath Contributors, supplied problems.json, record 3059 (OPG-148), OpenGarden collection. Catalogue identity and source transcription; CC BY 4.0.
\url{https://huggingface.co/datasets/ulamai/UnsolvedMath}.
\item[S2] Open Problem Garden, A nowhere-zero point in a linear mapping. Original problem and discussion; the mathematical formulation below is expanded from the supplied source record.
\url{http://www.openproblemgarden.org/op/a_nowhere_zero_point_in_a_linear_mapping}.
\item[S3] N. Alon and M. Tarsi, \emph{A nowhere-zero point in linear mappings}, Combinatorica 9 (1989), 393--395. Author-hosted original paper, checked 1 October 2026.
\url{https://web.math.princeton.edu/~nalon/PDFS/Publications/A%20nowhere-zero%20point%20in%20linear%20mappings.pdf}.
\item[C1] Reproducible finite-check code, executed 1 October 2026 with Python 3 (standard library only):
\path{scripts/check_attempt_batch_geometry_2026_10_01.py}.
\end{itemize}
\end{document}
